\chapter{Иллюстративный пример} \label{ch4}

В данной главе приведено пошаговое выполнение алгоритма GA-FP+Div на небольшом примере. 
Рассматривается тестовый набор из 5 тестовых случаев и 4 дефектов.

\section{Исходные данные} \label{ch4:sec1}

Пусть дан тестовый набор $T = \{t_1, t_2, t_3, t_4, t_5\}$ и множество дефектов 
$F = \{f_1, f_2, f_3, f_4\}$. Матрица покрытия дефектов и значения вероятности дефектов приведены в 
таблице~\ref{tab:example}.

\begin{table}[h]
\caption{Исходные данные иллюстративного примера}
\label{tab:example}
\begin{tabularx}{\textwidth}{|X|c|c|c|c|c|}
\hline
Тест & $f_1$ & $f_2$ & $f_3$ & $f_4$ & $r_i$ \\
\hline
$t_1$ & 1 & 0 & 1 & 0 & 0{,}3 \\
$t_2$ & 0 & 1 & 0 & 1 & 0{,}6 \\
$t_3$ & 1 & 1 & 0 & 0 & 0{,}4 \\
$t_4$ & 0 & 0 & 1 & 1 & 0{,}7 \\
$t_5$ & 1 & 0 & 0 & 1 & 0{,}5 \\
\hline
\end{tabularx}
\end{table}

Параметры алгоритма: $N = 2$ особи, $G = 1$ поколение, 
$p_c = 0{,}8$, $p_m = 0{,}2$, $w_c = 0{,}4$, $w_r = 0{,}4$, 
$w_d = 0{,}2$, $k = 3$.

\section{Шаг 1. Инициализация} \label{ch4:sec2}

Начальная популяция формируется случайным перемешиванием:
$$P_1 = [t_3, t_1, t_5, t_2, t_4], \quad P_2 = [t_2, t_4, t_1, t_5, t_3]$$

\section{Шаг 2. Вычисление фитнесса} \label{ch4:sec3}

Вычислим фитнесс для $P_1 = [t_3, t_1, t_5, t_2, t_4]$ при $k = 3$.

\textbf{Покрытие:} объединение покрытий всех тестов:
$$M_3 \cup M_1 \cup M_5 \cup M_2 \cup M_4 = \{f_1, f_2, f_3, f_4\}$$
$$\text{Cov}(P_1) = \frac{4}{4} = 1{,}0$$

\textbf{Вероятность дефектов} первых $k=3$ тестов 
($t_3$, $t_1$, $t_5$):
$$\text{FP}(P_1) = \frac{r_3 + r_1 + r_5}{3} = 
\frac{0{,}4 + 0{,}3 + 0{,}5}{3} = 0{,}400$$

\textbf{Структурное разнообразие} между соседними парами в топ-3:

Расстояние Жаккара между $t_3$ и $t_1$:
$$d_J(M_3, M_1) = 1 - \frac{|\{f_1\}|}{|\{f_1, f_2, f_3\}|} = 
1 - \frac{1}{3} = 0{,}667$$

Расстояние Жаккара между $t_1$ и $t_5$:
$$d_J(M_1, M_5) = 1 - \frac{|\{f_1\}|}{|\{f_1, f_3, f_4\}|} = 
1 - \frac{1}{3} = 0{,}667$$

$$\text{Div}(P_1) = \frac{0{,}667 + 0{,}667}{2} = 0{,}667$$

\textbf{Итоговый фитнесс:}
$$f(P_1) = 0{,}4 \cdot 1{,}0 + 0{,}4 \cdot 0{,}400 + 
0{,}2 \cdot 0{,}667 = 0{,}400 + 0{,}160 + 0{,}133 = 0{,}693$$

Аналогично для $P_2 = [t_2, t_4, t_1, t_5, t_3]$:
$$\text{FP}(P_2) = \frac{0{,}6 + 0{,}7 + 0{,}3}{3} = 0{,}533$$
$$d_J(M_2, M_4) = 1 - \frac{0}{4} = 1{,}000, \quad 
d_J(M_4, M_1) = 1 - \frac{1}{4} = 0{,}750$$
$$\text{Div}(P_2) = \frac{1{,}000 + 0{,}750}{2} = 0{,}875$$
$$f(P_2) = 0{,}4 \cdot 1{,}0 + 0{,}4 \cdot 0{,}533 + 
0{,}2 \cdot 0{,}875 = 0{,}400 + 0{,}213 + 0{,}175 = 0{,}788$$

Таким образом, $f(P_2) > f(P_1)$, поэтому $T^* = P_2$.

\section{Шаг 3. Кроссовер OX} \label{ch4:sec4}

Турнирная селекция выбирает $P_1$ и $P_2$ как родителей. 
Применяем OX-кроссовер с сегментом $[1, 3]$:

$$P_1 = [t_3,\ \underbrace{t_1,\ t_5,\ t_2}_{\text{сегмент}},\ t_4]$$
$$P_2 = [t_2,\ t_4,\ t_1,\ t_5,\ t_3]$$

Потомок $C_1$: копируем сегмент $[t_1, t_5, t_2]$ из $P_1$, 
остаток заполняем из $P_2$, пропуская $t_1$, $t_5$, $t_2$:
$$C_1 = [t_4,\ t_1,\ t_5,\ t_2,\ t_3]$$

\section{Шаг 4. Результат} \label{ch4:sec5}

Вычислим фитнесс потомка $C_1 = [t_4, t_1, t_5, t_2, t_3]$:
$$\text{FP}(C_1) = \frac{0{,}7 + 0{,}3 + 0{,}5}{3} = 0{,}500$$
$$d_J(M_4, M_1) = 0{,}750, \quad d_J(M_1, M_5) = 0{,}667$$
$$\text{Div}(C_1) = \frac{0{,}750 + 0{,}667}{2} = 0{,}708$$
$$f(C_1) = 0{,}4 \cdot 1{,}0 + 0{,}4 \cdot 0{,}500 + 
0{,}2 \cdot 0{,}708 = 0{,}400 + 0{,}200 + 0{,}142 = 0{,}742$$

Поскольку $f(C_1) = 0{,}742 < f(T^*) = 0{,}788$, лучшая особь 
не обновляется. Алгоритм возвращает:
$$T^* = [t_2, t_4, t_1, t_5, t_3]$$

Данная последовательность ставит вперёд тесты с высокой 
вероятностью дефектов ($t_2$: $r=0{,}6$, $t_4$: $r=0{,}7$) и 
разнообразным покрытием, что соответствует цели алгоритма.